% Part: sets-functions-relations
% Chapter: infinite
% Section: dedekinds-proof
%
\documentclass[../../../include/open-logic-section]{subfiles}

\begin{document}

\olfileid{sfr}{infinite}{dedekindsproof}

\olsection[``Bukti'' Dedekind]{``Bukti'' Dedekind Bab Anane
Himpunan Tanpa Wates}

Ing bab iki, kita wis menehi andharan wilangan natural kanthi teori
himpunan, yaiku nganggo aljabar Dedekind. Ing
\olref[arith][ref]{sec}, kita wis ngrembug teges filosofis saka
aritmetisasi analisis (lan bab liyane). Saiki kita kudu ngrembug teges
saka apa kang wis kasil ditindakake ing kene.

Ing saindhenging \olref[sfr][arith][]{chap}, kita nganggep wilangan
natural wis ana, banjur digunakake kanggo mbangun wilangan wutuh,
rasional, lan real kanthi cetha. Ing bab iki, kita ora menehi
pambangunan wilangan natural kanthi cetha. Kita mung nuduhake yen,
\emph{saka sembarang himpunan kang tanpa wates miturut Dedekind}, kita
bisa nemtokake himpunan kang tumindake persis kaya kang dikarepake
saka~$\Nat$.

Mula cetha yen kita ora bisa ngakoni wis mangsuli pitakon metafisika,
kayata \emph{obyek apa kang dadi wilangan natural}. Nanging iki malah
becik. Ing \olref[sfr][arith][ref]{sec}, kita negesake yen ora bener
nganggep definisi $\Real$ minangka himpunan potongan Dedekind dadi
\emph{panemon}, dudu katetepan kang migunani. Pengamatan kang wigati
yaiku potongan Dedekind menehi tuladha sipat matematis utama wilangan
real. Mangkono uga ing kene: pengamatan kang wigati yaiku \emph{saben}
aljabar Dedekind menehi tuladha sipat matematis utama wilangan natural.
(Pancen, Dedekind negesake bab iki kanthi mbuktekake yen kabeh aljabar
Dedekind iku \emph{isomorfis} (\citeyear[Teorema
132--3]{Dedekind1888}). Mula ora nggumunake yen akeh ``strukturalis''
saiki nyebut Dedekind minangka pelopor.)

Kajaba iku, kita wis nuduhake carane nyisipake teori wilangan natural
menyang teori himpunan prasaja kang naif. Teori iki dhewe isih rada
informal, nanging (katone) ora nganggep wilangan natural wis ana. Mula,
bisa uga kita lagi maju kanggo mujudake ancas Dedekind dhewe kang
ambisius, kang dijlentrehake kaya mangkene:
\begin{quote}
	Ing ngelmu, ora ana bab kang bisa dibuktekake kang kena dipracaya
	tanpa bukti. Senajan panjaluk iki katon lumrah, aku ora bisa nganggep
	yen panjaluk iki wis kawujud, malah ing cara-cara paling anyar kanggo
	nata dhasar ngelmu kang paling prasaja; yaiku perangan logika kang
	nangani teori wilangan. Nalika aku nyebut aritmetika (aljabar,
	analisis) mung minangka perangan logika, karepku yaiku konsep wilangan
	iku babar pisan ora gumantung marang gagasan utawa intuisi ngenani
	ruang lan wektu---malah, aku nganggep konsep iku minangka asil
	langsung saka hukum-hukum pamikiran murni.
	\citep[pambuka]{Dedekind1888}
\end{quote}
Gagasan Dedekind kang wani yaiku mangkene. Kita lagi wae nuduhake
carane mbangun wilangan natural mung nganggo teori himpunan (naif).
Ing \olref[sfr][arith][]{chap}, kita weruh carane mbangun wilangan real
saka wilangan natural lan sawatara teori himpunan. Mula, bisa uga
``aritmetika (aljabar, analisis)'' pungkasane ``mung dadi perangan
logika'' (miturut teges tembung ``logika'' kang dijembarakake Dedekind).

Mangkono gagasane. Nanging entenen sedhela. Pambangunan aljabar
Dedekind (sulih kita kanggo wilangan natural) gumantung marang anane
himpunan kang tanpa wates miturut Dedekind. (Delengen maneh
\olref[sfr][infinite][dedekind]{thm:DedekindInfiniteAlgebra}.) Yen
anane himpunan kuwi ora bisa ditetepake liwat ``logika'' utawa
``hukum-hukum pamikiran murni'', ancas iki mandheg.

Dadi, \emph{apa anane himpunan kang tanpa wates miturut Dedekind bisa
ditetepake nganggo ``hukum-hukum pamikiran murni''}? Iki upayane
Dedekind:
\begin{quote}
  Wewengkon pikiranku dhewe, yaiku totalitas $S$ saka kabeh bab kang
  bisa dadi obyek pikiranku, iku tanpa wates. Amarga yen $s$ nuduhake
  sawijine unsur saka~$S$, mula pikiran $s'$ yen~$s$ bisa dadi obyek
  pikiranku, iku dhewe uga unsur saka~$S$. Yen iki kita anggep minangka
  citra $\phi(s)$ saka unsur~$s$, mula \dots~$S$ iku tanpa wates
  [miturut Dedekind], kaya kang kudu dibuktekake.
	\citep[\S66]{Dedekind1888}
\end{quote}
Iki bab kang nggumunake banget amarga tinemu ing tengahing buku kang
bagean gedhene dumadi saka bukti matematis kang ketat banget. Ana rong
cathetan kang perlu diandharake.

Kapisan: ``bukti'' iki meh ora nduweni watak kang saiki bakal kita
arani ``matematis''. Bukti iki ngomongake obyek psikologis (pikiran),
lan malah pikiran kang mung \emph{bisa} ana.

Kapindho: paling ora manut cara kita ngandharake aljabar Dedekind,
``bukti'' iki duwe kekurangan teknis kang cetha. Yen argumen Dedekind
kasil, argumen iku mung netepake yen ana bab tanpa wates cacahe
(mligine, pikiran tanpa wates cacahe). Nanging Dedekind uga kudu menehi
alesan supaya~$S$ bisa dianggep minangka siji \emph{himpunan}, kang
duwe !!{element}s tanpa wates cacahe, dudu mung nganggep~$S$ minangka
\emph{sawetara bab} (jamak).

Kasunyatan yen Dedekind ora weruh lowongan iki bisa menehi panyana yen
panganggone tembung ``totalitas'' ora persis cocog karo panganggone
tembung ``himpunan'' \emph{kita}.\footnote{Pancen, ana alesan liyane
kanggo nganggep mangkono; delengen \citet[kaca~23]{Potter2004}.} Nanging
iki ora nggumunake banget. Ancas kang ditindakake sajrone rong bab
pungkasan---``pambangunan'' wilangan natural, banjur saka iku
``pambangunan'' wilangan wutuh, real, lan rasional---kabeh ditindakake
kanthi naif. Kita langsung nganggep himpunan iki utawa himpunan iku ana
nalika dibutuhake, tanpa netepake akeh asas umum bab persise himpunan
apa kang ana lan kapan anane. Nanging kita ngerti yen \emph{sawetara}
asas umum dibutuhake, amarga yen ora, kita bakal tiba ing Paradoks
Russell.

 Saiki wayahe kita ngluwihi kenaifan mau.

\end{document}
